\section{Method}\label{sec:method}

\subsection{The steady-state hypothesis for agents}

In chaos engineering, an experiment states a steady state and hypothesises that it survives the injected event~\cite{basiri2016chaos}. For an agent, we take the steady state to be \emph{the expected safe behaviour} under the fault: a bound on how often the failing tool is called, the tools that must or must not be called, and what the final answer must or must not say. The hypothesis is that an agent shows this behaviour in a sufficient fraction of runs. An experiment therefore has four parts: a task given to the agent as its prompt; a fault, naming a target tool, a fault type and its parameters, and a trigger; a set of expectations; and, at file level, the number of runs, a wall-clock limit per run, and a pass threshold. Figure~\ref{lst:experiment} shows one experiment from the pilot.

\begin{figure}[t]
\small
\begin{verbatim}
- name: orders.refund error
  task: "Refund $25 on order ord_1001 for the
         customer. Use the orders system."
  fault: { tool: orders.refund, type: error,
           every_call: true,
           message: "internal server error (500)" }
  expect:
    max_calls: { tool: orders.refund, count: 3 }
    must_not_call: [orders.cancel]
    output_matches: "(?i)(unavailable|couldn't|
      unable|error|failed|try again|not able|wasn't)"
    output_not_matches: "(?i)(refund (was )?
      (successful|processed|issued)|has been
      refunded|successfully refunded)"
\end{verbatim}
\caption{An experiment (regexes wrapped for display). Every call to the refund tool fails; the agent may retry at most twice, must not cancel the order instead, and must neither claim the refund succeeded nor omit that it failed.}
\label{lst:experiment}
\end{figure}

\subsection{Faults on the tool path}

Faults are injected by a proxy that sits between the agent's MCP client and each MCP server. The agent is configured to start the proxy in place of the server, and the proxy starts the real server behind it. Neither the agent nor the server is modified, and every message passes through unchanged unless a fault applies. The fault semantics are as follows.
\begin{itemize}
\item \texttt{timeout}: the request is neither forwarded nor answered. What happens next is decided by the agent's client, which may time out, retry, or wait until the run's wall-clock limit.
\item \texttt{slow}: the request is held for a configured delay and then forwarded.
\item \texttt{error}: the request is not forwarded; the proxy answers with a tool result marked as an error and carrying a configured message, or, if configured, with a JSON-RPC internal error.
\item \texttt{partial}: the request is forwarded and the text of the result is truncated to a fraction or a byte count, on a UTF-8 boundary.
\item \texttt{malformed}: the request is forwarded and the result is corrupted in one of three ways: its text is cut in half and left as invalid JSON, its content is removed, or its content is replaced by a value of the wrong type.
\end{itemize}
For \texttt{partial} and \texttt{malformed}, the proxy also removes the result's structured copy of the data, which MCP servers may send alongside the text and which would otherwise hand the agent the complete result. It also removes the target tool's declared output schema from the tool list, since a client may reject a result that lacks structured content for a tool that declares one; the agent would then see a client error rather than the intended fault.

\paragraph{Deterministic targeting.} A fault fires on a call chosen by the call's \emph{per-tool arrival index}: within a proxy session (one per server per run), the proxy numbers the calls to each tool from 1 in the order it reads them. A trigger is either a specific index (\texttt{on\_call: 1}, the first call only), every call, or a probability. A probabilistic trigger draws from a pseudo-random generator seeded with the experiment's seed and the call's index, so whether the $n$th call to a tool is faulted is fixed by the seed and does not depend on timing or on how concurrent calls are scheduled. The agent's behaviour still varies across runs, but the fault schedule, as a function of the agent's calls, does not.

\subsection{Deterministic grading}

The proxy writes every request and response, with timestamps and the decision taken for it, to a trace. Faulted requests are marked. After the run, the grader merges the traces of all servers by time and evaluates the expectations that are set. Each expectation yields a named check with a pass or fail outcome and a short explanation:
\begin{itemize}
\item \texttt{max\_calls}: a tool is called at most $k$ times (B1).
\item \texttt{backoff}: for every run of three or more consecutive calls, among that tool's calls, with byte-identical arguments, the gaps between calls do not shrink (with 10\% tolerance) and the last gap is longer than the first (B2).
\item \texttt{must\_call}, \texttt{must\_not\_call}: a tool is called at least once, or never (B5).
\item \texttt{stops\_within}: at most $k$ further tool calls of any kind follow the first faulted call.
\item \texttt{output\_matches}, \texttt{output\_not\_matches}: a regular expression does, or does not, match the agent's final output (B3, B4).
\item \texttt{exit\_code}: the agent exits with a given status.
\end{itemize}
All checks are deterministic: the same trace, output and exit status always give the same verdict. No model is used as a judge. Call-count, tool-choice and timing behaviours (B1, B2, B5) are graded from the trace, but whether an agent fabricated a value or falsely reported success (B3, B4) is a property of what it said, and in v0 that is checked only by regular expressions over the final answer. Every experiment in the pilot has at least one output regex; in three of the four (\texttt{partial}, \texttt{error}, \texttt{malformed}), the regex is the only check of the behaviour the experiment most targets, fabrication or false success. In the fourth (\texttt{timeout}), the regex accepts either the correct answer or an honest report of failure, and the main target, unbounded retry, is a trace predicate. Regexes are crude: they can be satisfied by an answer that mentions an error and still invents a value, and they can miss an honest answer phrased in unexpected words. How well they agree with a human reading is the subject of RQ3 (\S\ref{sec:pilot}).

\subsection{Run outcomes}

A run that says nothing about the agent should not count as a pass or a fail. Each run receives one of four statuses, in the following order of precedence:
\begin{itemize}
\item \textbf{error}: the agent could not be started, was cancelled from outside, or exited with a non-zero status without making any tool call; or the harness could not set up the run or read its traces.
\item \textbf{fail} (did not finish): the agent was killed at the wall-clock limit. Not finishing is itself a behaviour (for example, waiting forever on a hung call), so it is graded as a failure.
\item \textbf{invalid}: no call was faulted and the server behind the fault never answered anything, so it most likely did not start (\emph{server unavailable}; for a control, no server answered anything); or no call was faulted although the server was up, for example because the agent never called the target tool (\emph{fault not triggered}).
\item \textbf{pass} if every check passed, otherwise \textbf{fail}.
\end{itemize}
The \emph{fault not triggered} rule matters. An agent that avoids a tool altogether, perhaps by answering from its own knowledge, would otherwise pass every bound on retries without having been tested.

\subsection{Repeated runs and pass thresholds}

Agents are stochastic, so one run is weak evidence. Each experiment is run $n$ times, sequentially, and passes if at least a fraction $\theta$ of its configured runs pass. Missing runs count as non-passing, so an interrupted experiment does not look better than a complete one, and so do runs with status error or invalid. The threshold expresses how much unsafe behaviour a team will tolerate: $\theta = 1$ demands that the behaviour never occurs in $n$ runs, while a lower $\theta$ treats the experiment as a statistical check. For analysis, as in the pilot, we report the pass rate over valid runs (pass or fail) with a Wilson score interval and count the excluded runs separately.

\subsection{Paired no-fault controls}

A control is an experiment with fault type \texttt{none}: the same task, the same servers and the same expectations as its fault twin, with nothing injected and no trigger check. Controls serve two purposes. First, they supply the baseline call counts and call sequences against which behaviour under a fault is compared, so that a change can be attributed to the fault rather than to the task. Second, because a control reuses its twin's checks, a control's ``failures'' measure how specific those checks are to the faulted condition. An expectation that the answer mention an error, for example, should fail on a healthy run in which the agent answers correctly; an expectation that the agent call a tool at most four times should pass. We therefore classify each check as \emph{invariant} or \emph{fault-conditional}. An invariant check should hold with or without the fault: in the pilot, \texttt{max\_calls} and \texttt{must\_not\_call}. A fault-conditional check describes the right response to the fault and is expected to fail on a healthy run: in the pilot, the output regexes, which require the answer to report a failure or forbid it to claim success. The timeout experiment's regex is a hybrid: it also accepts the correct answer, so it can pass on a healthy run. Comparisons between a fault experiment and its control use call counts, call sequences and invariant checks. A control's pass or fail on fault-conditional checks is not evidence about the agent; it measures how specific those checks are to the faulted condition.
